\documentclass[runningheads]{llncs}

\usepackage{amsmath}
\usepackage{lipsum} 
\usepackage{booktabs} % For pretty tables
\usepackage{caption} % For caption spacing
\usepackage{subcaption} % For sub-figures
\usepackage{graphicx}
\usepackage{pgfplots}
\usepackage[all]{nowidow}
\usepackage[utf8]{inputenc}
\usepackage{tikz}
\usetikzlibrary{er,positioning,bayesnet}
\usepackage{multicol}
\usepackage{algpseudocode,algorithm,algorithmicx}
\usepackage{minted}
\usepackage{hyperref}
\usepackage[inline]{enumitem} % Horizontal lists

\newcommand{\card}[1]{\left\vert{#1}\right\vert}
\newcommand*\Let[2]{\State #1 $\gets$ #2}
\definecolor{blue}{HTML}{1F77B4}
\definecolor{orange}{HTML}{FF7F0E}
\definecolor{green}{HTML}{2CA02C}

\pgfplotsset{compat=1.14}

\renewcommand{\topfraction}{0.85}
\renewcommand{\bottomfraction}{0.85}
\renewcommand{\textfraction}{0.15}
\renewcommand{\floatpagefraction}{0.8}
\renewcommand{\textfraction}{0.1}
\setlength{\floatsep}{3pt plus 1pt minus 1pt}
\setlength{\textfloatsep}{3pt plus 1pt minus 1pt}
\setlength{\intextsep}{3pt plus 1pt minus 1pt}
\setlength{\abovecaptionskip}{2pt plus 1pt minus 1pt}

\begin{document}

% TITLE 
\title{TITLE}

% AUTHOR LIST
\author{Author1\inst{1} \and
Author2\inst{2}}

% INSTITUTE
\institute{University of Michigan, Ann Arbor, MI}

% MAKE THE TITLE ON THE PAGE
\maketitle   

% ABSTRACT
% \begin{abstract}
% your abstract here
% \end{abstract}

% PAPER SECTIONS

\section{Introduction}
You can make in text citations for anything in the bibliography (at the bottom of the document) like this \cite{Turing}. 


\section{Methods}


\textbf{Example Equation}. Below is Equation \ref{eq1}:

\begin{equation}\label{eq1}
    \mathbf{A} = \mathbf{U \Sigma V^{\text{T}}}
\end{equation}

\noindent \textbf{Example Multi-line Equation}. An aligned multi-line equation is shown in Equation \ref{eq2}:

\begin{align}\label{eq2}
    \mathbf{\dot{x}} &= \mu (\mathbf{x} - \frac{1}{3}\mathbf{x}^3 - \mathbf{y} ) \\
    \mathbf{\dot{y}} &= \frac{1}{\mu} \mathbf{x} 
\end{align}

\noindent You can also embed notation in the text like this: $\mathbf{Ax} = \lamda \mathbf{x}$.

\clearpage % forces this section to the next page
\section{Results}

\subsection{Example Subsection: Figures}
You can add figures to your document, as shown in Figure \ref{fig1}.

\begin{figure}[!tbh]
    \centering
    \includegraphics[scale=0.4]{figures/SVD_figure.png}
    \caption{\textbf{Example Figure.} Some other caption stuff here.}
    \label{fig1}
\end{figure}


\subsection{Example Subsection: Tables}

\bigskip % add some space after the figure
\noindent Tables may also be defined:

\begin{table}[!tbh]
    \centering
    \begin{tabular}{c|c}
         Input & $\mu x + 1$  \\
         \hline
         1 & 3 \\ 
         2 & 5 \\ 
         3 & 7 \\ 
         4 & 9  
    \end{tabular}
    \caption{\textbf{Example Table.} Other caption stuff.}
    \label{tab1}
\end{table}


\section{Conclusion}


% REFERENCES
\begin{thebibliography}{}
\bibitem{Turing}
A. M. Turing, “The chemical basis of morphogenesis,” Bltn Mathcal Biology, vol. 52, no. 1, pp. 153–197, Jan. 1990, doi: 10.1007/BF02459572.


\end{thebibliography}
\end{document}
